\documentclass[letterpaper,11pt]{article}
\usepackage{graphicx}
\usepackage{geometry}
\linespread{1.2}                        
\setlength{\oddsidemargin}{0.0in}            
\setlength{\textwidth}{6.5in}    
\usepackage[utf8]{inputenc} % Asegúrate de usar utf8
\usepackage[T1]{fontenc} % Asegúrate de usar T1 font encoding      
\setlength{\topmargin}{-.6in}            
\setlength{\textheight}{9in}
\usepackage{latexsym} 
\usepackage{amssymb} 
\usepackage{amsmath}
\usepackage{amsxtra}
\usepackage[mathcal]{eucal} 
\usepackage{enumerate}
\usepackage{hyperref}
\usepackage{rotating}
\usepackage{caption}
\usepackage{lscape}
\usepackage{paralist} 
\usepackage{indentfirst}
\usepackage[dvipsnames,svgnames]{xcolor} 
\usepackage{setspace}
\usepackage{multicol} 
\usepackage{float}

\definecolor{codegreen}{rgb}{0,0.6,0}
\definecolor{codegray}{rgb}{0.5,0.5,0.5}
\definecolor{codepurple}{cmyk}{0.65, 1, 0, 0.35}
\definecolor{backcolour}{rgb}{0.95,0.95,0.92}

\usepackage{listings}
\usepackage{xcolor}

\lstset{
    language=R, % lenguaje de programación
    basicstyle=\small\ttfamily, % estilo básico
    commentstyle=\color{gray}, % estilo para los comentarios
    keywordstyle=\color{blue}, % estilo para las palabras clave
    stringstyle=\color{black}, % estilo para las cadenas de texto
    showstringspaces=false, % no mostrar espacios en cadenas de texto
    tabsize=4, % tamaño de la tabulación
    breaklines=true, % permitir el salto de línea automático
    postbreak=\mbox{\textcolor{red}{$\hookrightarrow$}\space}, % símbolo al final de una línea partida
    numbers=left, % mostrar números de línea
    numberstyle=\tiny\color{gray}, % estilo para los números de línea
    frame=single, % mostrar un borde alrededor del código
    backgroundcolor=\color{gray!10}, % color de fondo
}

\begin{document}

\begin{center}
    {\Large  Taller 2: Econometría }

Profesor:    Pedro Comte Hidalgo\\
\end{center}

\textbf{Instrucciones:} 

\begin{itemize}
    \item Usted dispone hasta las 9:00 del día 28-08-2025.
    \item Cuide la presentación y redacción de sus respuestas.
    \item Puede utilizar su computador y los apuntes tanto de clase como de ayudantía.
    \item Debe enviar un archivo en formato PDF y un R script (extensión .R)
\end{itemize}

\section{Trabajo Base de Datos (30 puntos)}

Para responder las preguntas siguientes, utilice la base de datos \textit{ceosal1} de la librería \textbf{wooldridge}. Asegúrese de cargarla correctamente (sin nulos) y de especificar el programa de forma rigurosa para evitar cambios por aleatoriedad.

\begin{enumerate}

\item (10 puntos) Examine la base y elabore una breve descripción de sus variables, distinguiendo entre dependientes, explicativas, continuas y \textit{dummies}. \textquestiondown Para qué tipo de estudio económico se emplea habitualmente esta base de datos?



    \item (5 puntos) Genere una variable indicadora (\textit{dummy}) que tome el valor 1 cuando la variable \textit{sales} supere su promedio y 0 en caso contrario. Asigne a esta nueva variable un nombre que sea fácilmente identificable.
    


    \item (5 puntos) Cree una variable de la multuplicación entre la variable \textit{dummy} antes creada y la variable \textit{indus}. \textquestiondown Qué indica esta interacción?




    \item (5 puntos) Cree una variable de la multiplicación entre la variable continua \textit{roe}  y la variable \textit{dummy} antes creada. \textquestiondown Qué indica esta interacción?


\end{enumerate}

\section{Modelos de Regresión Lineal (70 puntos)}

\begin{enumerate}

    \item (15 puntos) Genere un modelo de regresión lineal, sobre alguna de las variables dependientes, con respecto a las variables \textit{sales}, \textit{roe} y \textit{ros}. Interprete sus resultados y la significancia de los regresores. \textquestiondown El modelo tiene significancia global al 5\%?



\item (10 puntos) Al modelo de regresión lineal anterior, agregue las variables creadas en el inciso 1.2 y 1.4, interprete sus resultados. \textquestiondown Cómo afectan los regresores de estas variables en el salario?



\item (15 puntos) Realiza un test de hipotesis para evaluar si las dos variables que se agregaron son significativas para el modelo simultaneamente. Use significancia del 5\%.


\item (15 puntos) Volviendo a la regresión del salario en función de \textit{sales}, \textit{roe} y \textit{ros}. Considere dos grupos, Grupo 1: Aquellos clientes que \textit{sales} es mayor al promedio, y Grupo 2: Aquellos clientes que \textit{sales} es menor al promedio. Realice un test de hipotesis para grupos al 5\% de significancia. Use el metodo que más le acomode. 

\item (15 puntos) Usando la regresión del salario en función de \textit{sales}, \textit{roe} y \textit{ros}, se distinguen los grupos mediante las \textit{dummies}; \textit{indus}, \textit{finance}, \textit{consprod} y \textit{utility}, que corresponden al tipo de empresa. Realice mediante test de Chow, con 5\% de significancia, si los modelos para cada grupo son iguales. 


\end{enumerate}





\end{document}
